% ================================================================
%  第 8 章  嵌入极小曲面的结构
% ================================================================
\chapter{嵌入极小曲面的结构}
\label{chap:8}

{\setlength{\emergencystretch}{2em}%
本章介绍关于 $\R^3$ 中嵌入极小曲面的若干新近结果. 第一节证明来自
\cite{CM4} 的一个局部结果: 一个嵌入极小圆盘要么是图, 要么在稍大的尺度上
包含一段\newterm{double-spiral-staircase}{双螺旋阶梯}{double spiral staircase}. 第二节陈述 \cite{CM6} 的\newterm{one-sided-curvature}{单侧曲率估计}{one-sided curvature estimate}. 第三节描述
\cite{CM10} 中单侧曲率估计的一个应用, 用以证明广义 Nitsche 猜想
(最初由 P. Collin \cite{Co} 证明). 之后我们转向 \cite{CM24} 对嵌入曲面的
Calabi--Yau 猜想的解决. 最后, 我们描述 \cite{CM7} 中关于有限亏格嵌入极小
曲面的主要结构定理, 以及若干依赖 \cite{CM3}--\cite{CM7} 结构理论的近期
唯一性结果.\par}

% ================================================================
\section{作为双螺旋阶梯的圆盘}
\label{sec:8.1}

我们先描述由 \cite{CM3}--\cite{CM6} 导出的真嵌入极小圆盘的局部分类.
这是理解有限亏格嵌入极小曲面的关键一步, 因为后者可以分解成若干圆盘或
``\newterm{pair-of-pants}{裤子形曲面}{pair of pants}''. 这一分解用到了引理~\ref{lem:1.5} 与引理~\ref{lem:1.11} 中的
拓扑单调性.

嵌入极小圆盘有两个经典模型. 第一个是定义在 $\R^2$ 中单连通区域
(例如平面本身)上的极小图, 第二个则像螺旋面那样的双螺旋阶梯. 一段
双螺旋阶梯由两个互相盘绕的阶梯组成, 使得两个人可以擦身而过而互不相遇.

在 \cite{CM3}--\cite{CM6} 中, 我们证明了这是仅有的可能情形; 事实上, 每个
嵌入极小圆盘要么是极小图, 要么可以由一段重新尺度化的螺旋面的一部分
逼近. 曲率小时它是图, 曲率超过某个阈值时它是螺旋面的一部分.

\cite{CM3}--\cite{CM6} 的结果, 对圆盘以及对其他拓扑类型都一样, 只用到极小
曲面的一个片段. 特别地, 这些曲面完全允许有边界; 例如, 当我们在下一个
定理中说``$\R^3$ 中任一嵌入极小圆盘要么是某个函数的图, 要么是双螺旋
阶梯的一部分''时, 我们的意思是: 若曲面包含在半径为 $r_0$ 的欧氏球内,
且边界包含在该球的边界上, 那么在一个半径为 $r_0$ 的某个固定(小)比例的
同心欧氏球内, 曲面的任一连通分支要么是某个函数的图, 要么是双螺旋阶梯
的一部分. 允许曲面有边界是一个要点, 它使这些结果格外有用. 注意, 由于
结论只针对``曲面的一个固定比例'', 这是一个内部估计.

下面是嵌入极小圆盘的主要结构定理\footnote{关于精确表述以及证明, 见
\cite{CM3}--\cite{CM6}.}:

\begin{theorem}
\label{thm:8.1}
(Colding--Minicozzi, \cite{CM3}--\cite{CM6}) $\R^3$ 中任一嵌入极小圆盘要么是
某个函数的图, 要么是双螺旋阶梯的一部分. 特别地, 若在某个点处曲率足够大,
则曲面是双螺旋阶梯的一部分(它可以由一段重新尺度化的螺旋面逼近).
另一方面, 若曲率处处低于某个阈值, 则曲面是某个函数的图.
\end{theorem}

定理~\ref{thm:8.1} 的证明有如下三个主要步骤:

\begin{enumerate}
\item[(A)] 固定一个整数 $N$(下文中曲率``大''的界将依赖于 $N$). 若一个嵌入
极小圆盘 $\Sigma$ 不是图(等价地, 若曲率在某个点处很大), 则它包含一个
$N$ 值的极小图, 而这一事实起初是在 $1/\max|A|$ 的尺度上建立的. 也就是说,
最初只证明该 $N$ 值图定义在一个圆环上, 其内外半径都与 $\max|A|$ 成反比.

\item[(B)] 这样一张位于 $\Sigma$ 内的、可能很小的 $N$ 值图, 随后可以看作
在 $\Sigma$ 内几乎一直延伸到边界, 即延伸为一个 $N$ 值图, 定义在一个外半径
与 $R$ 成正比的圆环上. 这里 $R$ 是 $\R^3$ 中包含 $\Sigma$ 边界的球的半径.

\item[(C)] 该 $N$ 值图不仅在水平方向(即与初始各片相切方向)延伸, 而且在
竖直方向(即与各片横截方向)也延伸. 也就是说, 一旦有了 $N$ 片, 就有了更多
的片; 事实上, 圆盘 $\Sigma$ 由两个多值图沿一条轴粘合而成.
\end{enumerate}

\mtfigure[0.52]{fig-8-1}{定理~\ref{thm:8.1} 的证明: A. 在 $\Sigma$ 中找到一张小的
$N$ 值图. B. 将其在 $\Sigma$ 中延伸为一张大的 $N$ 值图. C. 增加片数.
图片提供: Matthias Weber, www.indiana.edu/~minimal.}{fig:8.1}

证明超出本书范围, 但我们将给出 A 步的证明, 它依赖于第 2 章中我们证明过的
一些关于稳定与几乎稳定极小曲面的材料, 以及第 1 章中关于\newterm{multi-valued-graph}{多值图}{multi-valued graph}的结果.

A 步的准确表述是: 若一个球内的嵌入极小圆盘在某个点处曲率很大, 则它在
附近包含一张小的几乎平坦的多值图, 即:

\begin{theorem}
\label{thm:8.2}
(Colding--Minicozzi, \cite{CM4}) 给定 $N,\omega>1$ 及 $\epsilon>0$, 存在
$C=C(N,\omega,\epsilon)>0$ 使得下述成立:

设 $0\in\Sigma^2\subset B_R\subset R^3$ 是一个嵌入极小圆盘, 且
$\partial\Sigma\subset\partial B_R$. 若对某个 $0<r_0<R$ 有
\[
\sup_{B_{r_0}\cap\Sigma}|A|^2\leq 4|A|^2(0)=4C^2r_0^{-2},
\]
则存在 $\bar R<r_0/\omega$ 以及(在旋转 $\R^3$ 之后)一张 $N$ 值图
$\Sigma_g\subset\Sigma$, 它定义在 $D_{\omega\bar R}\setminus D_{\bar R}$ 上,
梯度 $\leq\epsilon$, 且 $\mathrm{dist}_\Sigma(0,\Sigma_g)\leq 4\bar R$.
\end{theorem}

回忆: 我们称 $N$ 值图 $\Sigma$ 的\emph{中间片} $\Sigma^M$ 是定义在
\begin{equation}
\bigl\{(\rho,\theta)\in\mathcal P\ \big|\ r_1<\rho<r_2
\text{ 且 } 0\leq\theta\leq 2\pi\bigr\}
\label{eq:8.1}
\end{equation}
之上的那部分.

\mtpara{嵌入极小圆盘的总曲率与面积}
利用第 2 章引理~\ref{lem:2.32} 的分解, 我们接下来得到: 对曲率有界的嵌入
极小圆盘, 其内蕴球的总曲率与面积有多项式型估计.

\begin{lemma}
\label{lem:8.3}
存在 $C_1$ 使得: 若 $0\in\Sigma\subset B_{2R}$ 是嵌入极小圆盘,
$\partial\Sigma\subset\partial B_{2R}$ 且 $|A|^2\leq 4$, 则
\begin{equation}
\int_0^R\int_0^t\int_{B_s^\Sigma}|A|^2\,ds\,dt
=2\bigl(\mathrm{Area}(B_R^\Sigma)-\pi R^2\bigr)
\leq 6\pi R^2+20C_1R^5.
\label{eq:8.2}
\end{equation}
\end{lemma}

\begin{proof}
设常数 $C_1$、截断函数 $\chi$ 以及集合 $\bigcup_j\Omega_j$ 由引理
\ref{lem:2.32} 给出, 使得每个 $\Omega_j$ 都是 $\frac12$-稳定的. 在
$B_R^\Sigma$ 上定义
\[
\psi=\psi\bigl(\mathrm{dist}_\Sigma(0,\cdot)\bigr)
=1-\mathrm{dist}_\Sigma(0,\cdot)/R,
\]
使得 $\chi\psi$ 在 $\frac12$-稳定子域的并 $\bigcup_j\Omega_j$ 之外为零.
把 $\chi\psi$ 代入 $\frac12$-稳定性不等式(见 \eqref{eq:2.136}), 并用吸收
不等式与 \eqref{eq:2.146}, 得
\begin{equation}
\begin{aligned}
\int|A|^2\chi^2\psi^2
&\leq 2\int|\nabla(\chi\psi)|^2\\
&=2\int\bigl(\chi^2|\nabla\psi|^2
+2\chi\psi\langle\nabla\chi,\nabla\psi\rangle
+\psi^2|\nabla\chi|^2\bigr)\\
&\leq 6C_1R^3+3\int\chi^2|\nabla\psi|^2
\leq 6C_1R^3+3R^{-2}\mathrm{Area}(B_R^\Sigma).
\label{eq:8.3}
\end{aligned}
\end{equation}
再用 \eqref{eq:2.145} 与 $|A|^2\leq 4$, 得
\begin{equation}
\int|A|^2\psi^2\leq 4C_1R^3+\int|A|^2\chi^2\psi^2
\leq 10C_1R^3+3R^{-2}\mathrm{Area}(B_R^\Sigma).
\label{eq:8.4}
\end{equation}
引理由 \eqref{eq:8.4} 与联系圆盘面积和总曲率的推论~\ref{cor:2.7} 得出.
\end{proof}

引理~\ref{lem:8.3} 中证明的多项式增长使我们能在两个不同的尺度上找到
具有固定加倍界的大内蕴球:

\begin{corollary}
\label{cor:8.4}
存在 $C_2$ 使得: 给定 $\beta,R_0>1$, 我们可以取 $R_2$ 使下述成立:

若 $0\in\Sigma\subset B_{R_2}$ 是嵌入极小圆盘, $\partial\Sigma\subset
\partial B_{R_2}$ 且
\[
\sup_\Sigma|A|^2\leq 4|A|^2(0)=4,
\]
则存在 $R_0<R<R_2/(2\beta)$ 使得
\begin{equation}
\int_{B_{3R}^\Sigma}|A|^2
+\beta^{-10}\int_{B_{2\beta R}^\Sigma}|A|^2
\leq C_2R^{-2}\mathrm{Area}(B_R^\Sigma).
\label{eq:8.5}
\end{equation}
\end{corollary}

\begin{proof}
设 $\mathcal A(s)=\mathrm{Area}(B_s^\Sigma)$. 给定 $m$, 由引理~\ref{lem:8.3} 得
\begin{equation}
\left(\min_{1\leq n\leq m}
\frac{\mathcal A((4\beta)^{2n}R_0)}{\mathcal A((4\beta)^{2n-2}R_0)}\right)^m
\leq\frac{\mathcal A((4\beta)^{2m}R_0)}{\mathcal A(R_0)}
\leq C_1'(4\beta)^{10m}R_0^3.
\label{eq:8.6}
\end{equation}
取 $m$ 使 $C_1'R_0^3<2^m$, 并令 $R_2=2(4\beta)^{2m}R_0$. 由
\eqref{eq:8.6}, 存在 $R_1=(4\beta)^{2n-2}R_0$($1\leq n\leq m$)使得
\begin{equation}
\frac{\mathcal A((4\beta)^2R_1)}{\mathcal A(R_1)}
\leq 2(4\beta)^{10}.
\label{eq:8.7}
\end{equation}
为简单起见, 设 $\beta=4^q$, $q\in\mathbf Z^+$. 与 \eqref{eq:8.6}、
\eqref{eq:8.7} 一样, 我们得到存在 $0\leq j\leq q$ 使
\begin{equation}
\frac{\mathcal A(4^{j+1}R_1)}{\mathcal A(4^jR_1)}
\leq\left[\frac{\mathcal A(4\beta R_1)}{\mathcal A(R_1)}\right]^{1/(q+1)}
\leq 2^{1/(q+1)}4^{10}.
\label{eq:8.8}
\end{equation}
令 $R=4^jR_1$. 结合 \eqref{eq:8.7}、\eqref{eq:8.8} 与推论~\ref{cor:2.7}
即得 \eqref{eq:8.5}. $\square$
\end{proof}



\mtpara{内蕴扇区}
本节引入\newterm{intrinsic-sector}{内蕴扇区}{intrinsic sector}的概念. 接下来我们证明: 足够大的``内蕴扇区''若几乎稳定, 则必含多值图. 这里用到的
``几乎稳定''的概念是 \cite{CM4} 第 2 节中的 $\delta$-稳定
(见定义~\ref{def:2.27}).

几乎稳定性有两个后果将被用到. 第一个是 $\delta$-稳定性不等式
\eqref{eq:2.136}. 第二个是下面这个关于 $\frac12$-稳定曲面的曲率估计, 它
由 \cite{CM2} 的论证给出(参见定理~\ref{thm:2.10}):

\begin{lemma}
\label{lem:8.5}
存在 $C$ 使得: 若 $\Sigma\subset\mathbb R^3$ 是双侧的 $\frac12$-稳定极小
曲面且 $B_R^\Sigma(x)\cap\partial\Sigma=\emptyset$, 则
\[
|A|^2(x)\leq\frac{C}{R^2}.
\]
\end{lemma}

\begin{proof}
这由定理~\ref{thm:2.10} 的证明稍作平凡修改即可得到. 见 \cite{CM2}.
\end{proof}

我们需要``内蕴扇区''的概念, 它定义在内蕴球边界的一条子曲线上. 这些扇区
使我们能够把圆环细分成若干子集, 类似于在极坐标中按 $\theta$ 坐标把欧氏
圆环细分. 具体地, 给定 $\gamma\subset\partial B_r^\Sigma$, 定义内蕴扇区
(见图~\ref{fig:8.2})
\begin{equation}
S_R(\gamma)=\left\{\exp_0(v)\ \big|\ r\leq|v|\leq r+R
\text{ 且 }\exp_0(rv/|v|)\in\gamma\right\}.
\label{eq:8.9}
\end{equation}
给定 $R>R_1>r$, 方便的记号是
\begin{equation}
S_{R_1,R}(\gamma)=S_R(\gamma)\setminus S_{R_1}(\gamma).
\label{eq:8.10}
\end{equation}

我们需要关于这些内蕴扇区的若干几何观察, 它们来自对 Riccati 方程的分析:

\mtfigure[0.46]{fig-8-2}{由 \eqref{eq:8.9} 定义的、位于曲线
$\gamma\subset\partial B_r^\Sigma$ 之上的内蕴扇区 $S_R(\gamma)$.
图片提供: Matthias Weber, www.indiana.edu/~minimal.}{fig:8.2}

\begin{lemma}
\label{lem:8.6}
设 $\Sigma$ 是包含 $0$ 的嵌入极小圆盘, $\gamma\subset\partial B_r^\Sigma$ 是一条
紧连通曲线, 且
\[
B_{r+R}^\Sigma\cap\partial\Sigma=\emptyset.
\]
定义法向指数映射
\[
\Phi:[0,R]\times\gamma\to S_R(\gamma)
\]
为 $\Phi(\rho,x)=\gamma_x(\rho)$, 其中 $\gamma_x$ 是从 $x$ 出发、与
$\partial B_r^\Sigma$ 法向(向外)的单位速测地线. 则:

\begin{enumerate}
\item[(1)] $\Phi$ 是一个微分同胚, $\Phi(\rho,\cdot):\gamma\to
S_R(\gamma)\cap\partial B_{r+\rho}^\Sigma$ 也是.

\item[(2)] 若给柱面 $[0,R]\times\gamma$ 赋以翘曲积(warped product)度量
$d\rho^2+(\rho/r+1)^2dx^2$, 则 $\Phi$ 不减小距离.

\item[(3)] 若 $k_g$ 是 $\rho$ 的等值集的高斯曲率, 则 $k_g\geq\frac{1}{\rho+r}$.

\item[(4)] 若 $\sup_{\{\rho\in(3s/4,s)\}}|A|^2\bar\rho^2\leq C_A$, 则
\[
\sup_{\{\rho=s\}}k_g\leq\frac{C_A+8}{s}.
\]

\item[(5)] 积分 $\int_{\{\rho=t\}}k_g$ 是 $t$ 的单调不减函数.
\end{enumerate}
\end{lemma}

\begin{proof}
由于 $\Sigma$ 是非正曲率的圆盘, 且 $B_{r+R}^\Sigma\cap\partial\Sigma=\emptyset$,
(标准的)指数映射 $\exp_0$ 是从 $B_{r+R}\subset R^2$ 到 $B_{r+R}^\Sigma$ 的
微分同胚. Gauss 引理表明 $\Sigma$ 中从 $0$ 出发的测地线与以 $0$ 为中心的
内蕴距离球正交. 这蕴含: 若 $v\in\mathbb R^2$ 满足 $r\leq|v|\leq r+R$ 且
$\exp_0(rv/|v|)\in\gamma$, 则
\begin{equation}
\exp_0(v)=\Phi\bigl(|v|-r,\exp_0(rv/|v|)\bigr).
\label{eq:8.11}
\end{equation}
由此 $\Phi$ 也是微分同胚, 得 (1); 而非正曲率的标准比较定理给出 (2) 与 (3).

为证 (4), 固定一条单位速测地线
\[
\sigma:[0,r+R]\to\Sigma
\]
满足 $\sigma(0)=0$, 并置 $f(t)=k_g\circ\sigma(t)$. 利用指数映射是微分同胚
且 $\Sigma$ 是曲面, 一个标准计算给出
\begin{equation}
f'(t)=-f^2(t)-K_\Sigma\circ\sigma(t)
=-f^2(t)+\frac12|A|^2\circ\sigma(t),
\label{eq:8.12}
\end{equation}
其中最后一个等号用了 Gauss 方程与极小性. (4) 中的假设给出, 对
$t\in[r+3s/4,r+s]$ 有
\begin{equation}
f'(t)\leq-f^2(t)+\frac89C_As^{-2}<-f^2(t)+C_As^{-2}.
\label{eq:8.13}
\end{equation}
现在分两种情形讨论:

\textbf{情形 1:} 先设存在某个 $\bar t\in[r+3s/4,r+s]$ 使得
\begin{equation}
f(\bar t)<\frac{\sqrt{2C_A}}{s}.
\label{eq:8.14}
\end{equation}
由于 \eqref{eq:8.13} 蕴含在 $[\bar t,r+s]$ 上 $f'(t)\leq C_As^{-2}$,
微积分基本定理给出
\begin{equation}
f(r+s)\leq f(\bar t)+(r+s-\bar t)\frac{C_A}{s^2}
\leq\frac{\sqrt{2C_A}}{s}+\frac{C_A}{4s}
<\frac{C_A+1}{s},
\label{eq:8.15}
\end{equation}
其中最后一个不等式用了 $\sqrt{2C_A}=2\sqrt{C_A/2}\leq C_A/2+1$.

\textbf{情形 2:} 反之设对所有 $t\in[r+3s/4,r+s]$ 都有
$f(t)\geq\frac{\sqrt{2C_A}}{s}$. 由 \eqref{eq:8.13} 得, 在该区间上
\begin{equation}
f'\leq-f^2+\frac{C_A}{s^2}\leq-\frac12f^2.
\label{eq:8.16}
\end{equation}
由于 $f(t)\geq\frac1t$, $1/f$ 有意义, 由 \eqref{eq:8.16} 计算得
\begin{equation}
(1/f)'=-\frac{f'}{f^2}\geq\frac12,
\label{eq:8.17}
\end{equation}
于是微积分基本定理给出
\begin{equation}
\frac1{f(r+s)}\geq\frac1{f(r+3s/4)}+\frac{s}{8}\geq\frac{s}{8}.
\label{eq:8.18}
\end{equation}
我们得到 $f(r+s)\leq\frac8s$.

因此, 无论哪种情形都有 $f(r+s)\leq\frac{C_A+8}{s}$, (4) 随之成立.

最后一个性质 (5) 直接由带角 Gauss--Bonnet 定理(见 \cite{dC1} 第 274 页)、
Gauss 引理(它给出各角处的角)以及由 Gauss 方程与极小性带来的曲率非正性
得出.
\end{proof}

\mtpara{几乎稳定的扇区含有多值图}
在下一个引理中, 我们用 $\mathcal T_s(\Omega)$ 表示集合 $\Omega$ 的
$s$-管状邻域:
\[
\mathcal T_s(\Omega)=\{x\in\Sigma\ \big|\ \mathrm{dist}_\Sigma(x,\Omega)<s\}.
\]
下一个引理是 \cite{CM3} 中控制几乎稳定扇区面积增长的一个结果的小幅推广.
它仿照第 2 章中 \cite{CM2} 对稳定圆盘的结果.

\begin{lemma}
\label{lem:8.7}
(Colding--Minicozzi, \cite{CM3}) 存在 $C$ 使得: 若 $0\in\Sigma$ 是嵌入极小
圆盘, 含有一条曲线 $\gamma\subset\partial B_{r_1}^\Sigma$ 且

\begin{itemize}
\item $\int_\gamma k_g<C_0m$ 且 $\mathrm{Length}(\gamma)\leq3\pi m r_1$;

\item $\mathcal T_{r_1/8}(S_{R_1}(\gamma))$ 是 $\frac12$-稳定的且与
$\partial\Sigma$ 不交;
\end{itemize}

则对任意 $t\in(2r_1,3R_1/4)$ 有
\begin{equation}
\int_\gamma k_g
\leq\frac{\mathrm{Length}\bigl(\partial B_{r_1+t}^\Sigma\cap S_{R_1}(\gamma)\bigr)}{t}
\leq C\left((C_0+1)m+\frac{R_1}{r_1}\right).
\label{eq:8.19}
\end{equation}
此外, 对任意 $\Omega>2$ 有
\begin{equation}
\int_{S_{\Omega r_1,R_1/\Omega}(\gamma)}|A|^2
\leq C\left(\frac{R_1}{r_1}+\frac{(C_0+1)m}{\log\Omega}\right).
\label{eq:8.20}
\end{equation}
\end{lemma}

\begin{proof}
在本证明中, 为简洁起见我们把扇区 $S_t(\gamma)$ 记作 $S_t$. 定义
$S_{R_1}$ 上的函数 $\rho$ 为
\[
\rho(x)=\mathrm{dist}_{S_{R_1}}(x,\gamma).
\]
由引理~\ref{lem:8.6}, 这里的 $\rho$ 与那里的 $\rho$ 一致.

定义函数 $\ell(t)$ 与 $K(t)$ 如下:
\begin{equation}
\ell(t)=\mathrm{Length}\bigl(\{\rho=t\}\bigr),
\label{eq:8.21}
\end{equation}
\begin{equation}
K(t)=\int_{S_t}|A|^2.
\label{eq:8.22}
\end{equation}
由于指数映射是嵌入, 且垂直于 $\partial\gamma$ 的测地线与 $\rho$ 的等值集
正交(由 Gauss 引理), 弧长的一阶变分公式给出
\begin{equation}
\ell'(t)=\int_{\{\rho=t\}}k_g>0.
\label{eq:8.23}
\end{equation}
设 $d\mu$ 为 $\rho$ 的等值集上的一维 Hausdorff 测度. Jacobi 方程给出
\begin{equation}
\frac{d}{dt}(k_g\,d\mu)=-K_\Sigma\,d\mu=\frac{|A|^2}{2}\,d\mu,
\label{eq:8.24}
\end{equation}
其中第二个等号用了极小性与 Gauss 方程. 定义 $\bar K(t)$ 为 $K(t)$ 的积分:
\[
\bar K(t)=\int_0^tK(s)\,ds.
\]
把 \eqref{eq:8.24} 积分两次, 由 \eqref{eq:8.23} 得
\begin{equation}
\ell(t)=\ell(0)+\int_0^t\left(\int_\gamma k_g+\frac{K(s)}2\right)ds
=\mathrm{Length}(\gamma)+t\int_\gamma k_g+\frac{\bar K(t)}2.
\label{eq:8.25}
\end{equation}
由于 $\bar K(t)\geq0$, 这给出 \eqref{eq:8.19} 中的第一个不等式. 用余面积
公式与 $|\nabla\rho|=1$, 再积分 \eqref{eq:8.25}, 得
\begin{equation}
\begin{aligned}
\frac{\mathrm{Area}(S_{R_1})}{R_1^2}
&=\frac1{R_1^2}\int_0^{R_1}\ell(t)\\
&\leq\frac{\mathrm{Length}(\gamma)}{R_1}+\int_\gamma\frac{k_g}2
+\frac1{R_1^2}\int_0^{R_1}\frac{\bar K(t)}2\\
&\leq3\pi m+\frac12C_0m+\frac1{R_1^2}\int_0^{R_1}\frac{\bar K(t)}2.
\label{eq:8.26}
\end{aligned}
\end{equation}
其中最后一个不等式用了(假设的)$\gamma$ 的长度与总曲率界.

我们将用 $\frac12$-稳定性来控制总曲率, 但这引入了我们在第 2 章证明稳定
内蕴球的曲率估计时不需要处理的新项. 我们需要为``内边界'' $\gamma$、
``外边界'' $\{\rho=R_1\}$ 以及``侧边''(即垂直于 $\partial\gamma$ 的测地线)
准备截断函数.

在 $S_{R_1}$ 上定义函数 $\psi$ 为
\[
\psi=\psi(\rho)=1-\rho/R_1.
\]
接着, 令 $\gamma_a$ 与 $\gamma_b$ 是垂直于 $\partial\gamma$ 的测地线, 并置
\[
d_S=\mathrm{dist}_{S_{R_1}}(\cdot,\gamma_a\cup\gamma_b).
\]
在 $S_{R_1}$ 上定义函数 $\chi_1,\chi_2$ 为
\begin{equation}
\chi_1=\chi_1(d_S)=
\begin{cases}
d_S/r_1&\text{若 }0\leq d_S\leq r_1,\\
1&\text{其他},
\end{cases}
\label{eq:8.27}
\end{equation}
\begin{equation}
\chi_2=\chi_2(\rho)=
\begin{cases}
\rho/r_1&\text{若 }0\leq\rho\leq r_1,\\
1&\text{其他}.
\end{cases}
\label{eq:8.28}
\end{equation}
置 $\chi=\chi_1\chi_2$. 由于对每个 $x\in S_{R_1}$ 都有 $B_{r_1/8}^\Sigma(x)$
是 $\frac12$-稳定的, 引理~\ref{lem:8.7} 的曲率估计给出
\begin{equation}
\sup_{S_{R_1}}|A|^2\leq Cr_1^{-2}.
\label{eq:8.29}
\end{equation}
用它以及标准比较定理来估计 $\gamma\cup\gamma_a\cup\gamma_b$ 的管状邻域
面积, 得
\begin{equation}
\mathrm{Area}\bigl(S_{R_1}\cap\{\chi<1\}\bigr)
\leq C'\bigl(R_1r_1+mr_1^2\bigr).
\label{eq:8.30}
\end{equation}
把这个面积界与 \eqref{eq:8.29} 结合, 得
\begin{equation}
\int_{S_{R_1}}|\nabla\chi|^2
+\int_{S_{R_1}\cap\{\chi<1\}}|A|^2
\leq\tilde C\left(\frac{R_1}{r_1}+m\right).
\label{eq:8.31}
\end{equation}
在证明末尾, 我们还需要下面这个关于 $\chi_1$ 的界(我们正是用它得到
\eqref{eq:8.31}):
\begin{equation}
\int_{S_{R_1}}|\nabla\chi_1|^2
+\int_{S_{R_1}\cap\{\chi_1<1\}}|A|^2
\leq\tilde C\frac{R_1}{r_1}.
\label{eq:8.32}
\end{equation}
把 $\chi\psi$ 代入 $\frac12$-稳定性不等式, 并用吸收不等式
$2ab\leq\epsilon a^2+\frac1\epsilon b^2$ 与 \eqref{eq:8.31}, 得
\begin{equation}
\begin{aligned}
\frac12\int|A|^2\chi^2\psi^2
&\leq\int|\nabla(\chi\psi)|^2\\
&=\int\bigl(\chi^2|\nabla\psi|^2
+2\chi\psi\langle\nabla\chi,\nabla\psi\rangle
+\psi^2|\nabla\chi|^2\bigr)\\
&\leq\tilde C(1+1/\epsilon)\left(\frac{R_1}{r_1}+m\right)
+(1+\epsilon)\int\chi^2|\nabla\psi|^2\\
&\leq\tilde C(1+1/\epsilon)\left(\frac{R_1}{r_1}+m\right)
+\frac{(1+\epsilon)}{R_1^2}\int_0^{R_1}\ell(t)\,dt.
\label{eq:8.33}
\end{aligned}
\end{equation}
其中除最后一个积分外, 所有积分都在 $S_{R_1}$ 上取.

利用 $(\psi^2)''=2/R_1^2$, 再在一元情形下分部积分两次, 得
\begin{equation}
\begin{aligned}
\frac2{R_1^2}\int_0^{R_1}\bar K(t)
&=\int_0^{R_1}\bar K(t)(\psi^2)''\\
&=-\int_0^{R_1}K(t)(\psi^2)'
=\int_{S_{R_1}}|A|^2\psi^2\\
&\leq\tilde C\left(\frac{R_1}{r_1}+m\right)
+\int_{S_{R_1}}|A|^2\chi^2\psi^2\\
&\leq\tilde C(3+2/\epsilon)\left(\frac{R_1}{r_1}+m\right)
+2(1+\epsilon)R_1^{-2}\int_0^{R_1}\ell(t).
\label{eq:8.34}
\end{aligned}
\end{equation}
其中第一个不等式用了 \eqref{eq:8.31}, 最后一个不等式用了
\eqref{eq:8.33}.

(把 $\psi$ 同时看作 $[0,R_1]$ 上与 $S_{R_1}$ 上的函数, 略有记号上的滥用.
第三个等号用了余面积公式与 $|\nabla\rho|=1$ 来求 $K(t)$ 的导数.)

取 $\epsilon=\frac12$, 并结合 \eqref{eq:8.26} 与 \eqref{eq:8.34}, 得
\begin{equation}
\begin{aligned}
\frac4{R_1^2}\int_0^{R_1}\ell(t)
&\leq12\pi m+2C_0m+\frac2{R_1^2}\int_0^{R_1}\bar K(t)\\
&\leq(12\pi+2C_0)m+\frac3{R_1^2}\int_0^{R_1}\ell(t)
+7\tilde C\left(\frac{R_1}{r_1}+m\right).
\label{eq:8.35}
\end{aligned}
\end{equation}
关键在于 $\ell(t)$ 的积分在两边都出现, 但右边的常数更小, 于是得到一个界.
即 \eqref{eq:8.35} 给出
\begin{equation}
R_1^{-2}\mathrm{Area}(S_{R_1})
=R_1^{-2}\int_0^{R_1}\ell(t)
\leq C_4\left(\frac{R_1}{r_1}+m(1+C_0)\right).
\label{eq:8.36}
\end{equation}
由于 $\ell(t)$ 单调递增(由 \eqref{eq:8.23}), \eqref{eq:8.36} 给出
\eqref{eq:8.19} 中的第二个不等式(取 $t=3R_1/4$). 由于上面的论证在把
$R_1$ 换成满足 $2r_1<t<R_1$ 的 $t$ 时同样适用, 我们得到 \eqref{eq:8.19}
在 $2r_1\leq t\leq3R_1/4$ 上成立.

为证最后一个断言, 我们将用稳定性不等式以及对数截断技巧, 来利用面积的
二次增长. 定义截断函数 $\psi_1$ 为
\begin{equation}
\psi_1=\psi_1(\rho)=
\begin{cases}
\log(\rho/r_1)/\log\Omega&\text{在 }S_{r_1,\Omega r_1}\text{ 上},\\
1&\text{在 }S_{\Omega r_1,R_1/\Omega}\text{ 上},\\
-\log(\rho/R_1)/\log\Omega&\text{在 }S_{R_1/\Omega,R_1}\text{ 上},\\
0&\text{其他}.
\end{cases}
\label{eq:8.37}
\end{equation}
用 \eqref{eq:8.19} 与 \eqref{eq:8.36}(如 \eqref{eq:1.107} 中那样), 得
\begin{equation}
\int_{S_{R_1}}|\nabla\psi_1|^2
\leq\frac{C(C_0+m+R_1/r_1)}{\log\Omega}.
\label{eq:8.38}
\end{equation}
与 \eqref{eq:8.33} 一样, 把 $\frac12$-稳定性不等式用于 $\chi_1\psi_1$, 得
\begin{equation}
\begin{aligned}
\frac12\int|A|^2\chi_1^2\psi_1^2
&\leq\int_{S_{R_1}}|\nabla(\psi_1\chi_1)|^2\\
&\leq2\int_{S_{R_1}}\bigl(|\nabla\psi_1|^2+|\nabla\chi_1|^2\bigr)\\
&\leq\frac{2C((C_0+1)m+R_1/r_1)}{\log\Omega}
+\frac{2\tilde C R_1}{r_1}.
\label{eq:8.39}
\end{aligned}
\end{equation}
最后, 结合 \eqref{eq:8.32} 与 \eqref{eq:8.39} 得 \eqref{eq:8.20}.
\end{proof}


% ------------------------------------------------------------
%  第 8 章  §1.3 续 + §1.4
% ------------------------------------------------------------
来自 \cite{CM3} 的下一个结果表明, 一个大的几乎稳定的内蕴扇区必定含有一个
多值图. 其思路是: 前一个引理给出总曲率很小的子扇区, 因而是局部图形化的;
若这些子扇区足够``宽'', 它们就含有多值图.

\begin{proposition}[Colding--Minicozzi, \cite{CM3}]
\label{prop:8.8}
给定 $\omega>8$、$1>\epsilon>0$、$C_0$ 与 $N$, 存在 $m_1,\Omega_1$, 使得下述
结论成立:

若 $0\in\Sigma$ 是嵌入极小圆盘, 且含有一条曲线
$\gamma\subset\partial B_{r_1}^\Sigma$ 满足

\begin{itemize}
\item $\displaystyle\int_\gamma k_g<C_0m_1$ 且
$\operatorname{Length}(\gamma)=m_1r_1$;
\item $\mathcal{T}_{r_1/8}\bigl(S_{\Omega_1^2\omega r_1}(\gamma)\bigr)$ 是
$1/2$-稳定的且与 $\partial\Sigma$ 不交,
\end{itemize}

则(在旋转 $\R^3$ 之后)$S_{\Omega_1^2\omega r_1}(\gamma)$ 含有一个 $N$ 值图
$\Sigma_N$, 其定义在 $D_{\omega\Omega_1r_1}\setminus D_{\Omega_1r_1}$ 上, 梯度
不超过 $\epsilon$, $|A|\leq\epsilon/r$, 且
\begin{equation}
\operatorname{dist}_{S_{\Omega_1^2\omega r_1}(\gamma)}(\gamma,\Sigma_N)
 <4\Omega_1r_1.
\label{eq:8.40}
\end{equation}
\end{proposition}

\begin{proof}
与前一个证明一样, 定义函数 $\rho$ 为
\[
\rho(x)=\operatorname{dist}_{S_{\Omega_1^2\omega r_1}}(x,\gamma).
\]
我们将取 $\Omega_1>12$, 然后令 $m_1=\omega\Omega_1^2\log\Omega_1$. 由引理
\ref{lem:8.7}(其中 $\Omega=\Omega_1/6$、$R_1=\Omega_1^2\omega r_1$、
$m=32m_1/3$)得
\begin{equation}
\begin{aligned}
\int_{S_{\frac{\Omega_1r_1}{6},\,6\Omega_1\omega r_1}(\gamma)}|A|^2
&\leq C\left(\Omega_1^2\omega+\frac{m_1(1+C_0)}{\log(\Omega_1/6)}\right)\\
&\leq C\left(\Omega_1^2\omega+\frac{2m_1(1+C_0)}{\log\Omega_1}\right)\\
&=\frac{(3+2C_0)Cm_1}{\log\Omega_1}.
\end{aligned}
\label{eq:8.41}
\end{equation}
把 $\gamma$ 分成 $m_2\approx m_1/(8\pi(N+1))$ 条曲线
$\gamma_1,\ldots,\gamma_{m_2}\subset\gamma$, 使
\[
\operatorname{Length}(\gamma_i)=8\pi(N+1)r_1.
\]
由于我们已把 $\gamma$ 上的扇区分成内部两两不交的 $m_2$ 个子扇区, 至少有一个
$\gamma_i$ 满足
\begin{equation}
\begin{aligned}
\int_{S_{\frac{\Omega_1r_1}{6},\,6\Omega_1\omega r_1}(\gamma_i)}|A|^2
&\leq\frac1{m_2}
  \int_{S_{\frac{\Omega_1r_1}{6},\,6\Omega_1\omega r_1}(\gamma)}|A|^2\\
&\leq\frac{(3+2C_0)C\,8\pi(N+1)}{\log\Omega_1},
\end{aligned}
\label{eq:8.42}
\end{equation}
其中最后一个不等式用了 \eqref{eq:8.41}.

只要我们把 $\Omega_1$ 取成比 $N$ 大得多, \eqref{eq:8.42} 就给出一个非常平坦的
扇区. 我们将用这种平坦性来得到所求的 $N$ 值图.

引理~\ref{lem:8.6} 中的性质 (3) 蕴含
\[
\int_{\gamma_1/2}k_g\geq\frac1{r_1}\operatorname{Length}(\gamma_1/2)
 =4\pi(N+1),
\]
而引理~\ref{lem:8.6} 中的性质 (5) 蕴含
\[
\int_{S_{\Omega_1^2\omega r_1}(\gamma_i/2)\cap\{\rho=t\}}k_g
\]
关于 $t$ 单调不减. 把这两点结合起来, 我们可以选取一条连通曲线
$\tilde\gamma\subset\gamma_i/2$, 使
\begin{equation}
\int_{S_{\Omega_1^2\omega r_1}(\tilde\gamma)\cap
  \{\rho=\Omega_1r_1/3\}}k_g=2\pi(N+1).
\label{eq:8.43}
\end{equation}
置 $S=S_{\Omega_1r_1/3,\,3\Omega_1\omega r_1}(\tilde\gamma)$ 以及
$\hat\gamma=S\cap\{\rho=\Omega_1r_1/3\}$.

接下来, 我们用 Gauss--Bonnet 定理、\eqref{eq:8.42} 与 \eqref{eq:8.43}(当
$\Omega_1$ 足够大时)得到
\begin{equation}
2\pi(N+1)\leq\int_{S\cap\{\rho=t\}}k_g
\leq2\pi(N+1)+\frac12\int_S|A|^2\leq2\pi(N+2).
\label{eq:8.44}
\end{equation}
把总曲率的上界与引理~\ref{lem:8.6} 中逐点的下界
$k_g\geq1/(\rho+r_1)$ 结合起来, 得到长度界
\[
\operatorname{Length}(S\cap\{\rho=t\})\leq2\pi(N+2)(t+r_1)
\leq4\pi(N+2)t.
\]
最后注意, 由稳定性、\eqref{eq:8.42} 以及引理~\ref{lem:8.6} 中的 (2), 均值
定理给出: 对 $y\in S$,
\begin{equation}
\sup_{\mathcal{B}_{\rho(y)/3}(y)}|A|^2
\leq\frac{C_1N}{\rho^2(y)\log\Omega_1}.
\label{eq:8.45}
\end{equation}
沿 $\rho$ 的射线与水平集积分 \eqref{eq:8.45}, 得
\begin{equation}
\max_{x,y\in S}\operatorname{dist}_{\mathbf{S}^2}\bigl(N(x),N(y)\bigr)
\leq\frac{C_2(\log\omega+N)\sqrt{N}}{\sqrt{\log\Omega_1}}.
\label{eq:8.46}
\end{equation}
(其中的 $\log\omega$ 项来自沿射线的积分, 因为外半径与内半径之比约为 $\omega$;
另一项来自沿水平集 $\rho=t$ 的积分, 该水平集长度至多为 $4\pi(N+2)t$, 而
$|A|$ 至多约为 $\sqrt{N/\log\Omega_1}\,t^{-1}$.)

现在我们可以把这些事实结合起来得到本命题; 这与第 2 章引理~\ref{lem:2.18} 的
证明相当类似. 取 $\Omega_1$ 使得
\[
\frac{C_2(\log\omega+N)\sqrt{N}}{\sqrt{\log\Omega_1}}<C_3\epsilon,
\]
其中 $C_3$ 是稍后选取的一个小常数. 当 $C_3$ 充分小时, 在旋转 $\R^3$ 之后, $S$
在 $\{x_3=0\}$ 上局部是一个梯度不超过 $\epsilon/2$ 的图. 由于
$\tilde\gamma\subset B_{2r_1}$ 且 $\Omega_1>12$, 我们有
\[
\hat\gamma\subset B_{2r_1+\Omega_1r_1/3}\subset B_{\Omega_1r_1/2}.
\]
再把 $\Omega_1$ 取大一些, 并结合引理~\ref{lem:8.6} 中的 (3)、\eqref{eq:8.44}、
\eqref{eq:8.45} 与 \eqref{eq:8.46}, 我们看到(正交投影)$\Pi(\hat\gamma)$ 是一条
凸平面曲线, 其总曲率至少为 $(2N+1)\pi$, 因而其 Gauss 映射至少 $N$ 次覆盖
$S^1$. 给定 $x\in\tilde\gamma$, 置
$\tilde\gamma_x=S\cap\gamma_x$. 由 \eqref{eq:8.45}, $\tilde\gamma_x$ 的总
(外蕴测地)曲率至多为
\[
\frac{C_2\sqrt{N}\log\omega}{\sqrt{\log\Omega_1}}<C_3\epsilon,
\]
于是 $\tilde\gamma_x$ 落在一个以 $\tilde{x}=\tilde\gamma_x\cap\hat\gamma$ 处的
切向射线为中心的窄锥内. 当 $C_3$ 足够小时, 这蕴含 $\tilde\gamma_x$ 不发生
旋转, 并且
\begin{equation}
\bigl|\Pi(\tilde{x})-\Pi(\tilde\gamma_x\cap\{\rho=t\})\bigr|
\geq\frac9{10}\left(t-\frac{\Omega_1r_1}{3}\right).
\label{eq:8.47}
\end{equation}
因此, $\Pi\bigl(\partial\tilde\gamma_x\setminus\{\tilde{x}\}\bigr)\notin
D_{2\omega\Omega_1r_1}$, 这给出 $\Sigma_d$, 并且还给出
\[
\operatorname{dist}_{S_{\Omega_1^2\omega r_1}(\gamma)}(\gamma,\Sigma_d)
<2\Omega_1r_1.
\]
\end{proof}

\mtpara{轴附近的局部结构}
证明定理~\ref{thm:8.2} 的关键, 是找到 $n$ 个大的内蕴扇区, 并给出与尺度无关的
曲率界. 为此:

\begin{itemize}
\item 我们首先使用极小圆盘中内蕴球的曲率估计, 即定理~\ref{thm:2.17}, 从下方
控制 $\operatorname{Length}(\partial B_R^\Sigma)/R$, 其中 $R\geq R_0$.
\item 推论~\ref{cor:8.4} 给出 $R_3>R_0$ 以及 $n$ 条长的互不相交的曲线
$\tilde\gamma_i\subset\partial B_{R_3}^\Sigma$, 使得 $\tilde\gamma_i$ 上的扇区
有有界的 $\int|A|^2$.
\item 定理~\ref{thm:2.17} 给出与尺度无关的曲率界.
\item 一旦有了这些扇区, 当 $n$ 足够大时, 必有两个在 $\R^3$ 中很接近, 从而由
第 2 章的引理~\ref{lem:2.29} 与引理~\ref{lem:2.30}, 它们是 $1/2$-稳定的.
\item 最后, $N$ 值图由 \cite{CM3} 的推论 II.1.34 给出(见命题~\ref{prop:8.8}).
\end{itemize}

\mtfigure[0.44]{fig-8-3}{方程 \eqref{eq:8.51} 把一个穿孔球分成扇区
$\tilde S_i$.}{fig:8.3}

\begin{proof}[定理~\ref{thm:8.2} 的证明]
把 $\Sigma$ 按 $C/r_0$ 重新尺度化, 使得 $|A|^2(0)=1$ 且在 $B_C$ 上
$|A|^2\leq4$.

设 $C_2$ 由推论~\ref{cor:8.4} 给出, 再令 $m_1,\Omega_1>\pi$ 由命题
\ref{prop:8.8} 给出, 其中那里的 $C_0=2C_2+2$. 固定一个很大的 $a_0$
(稍后选取).

由于我们有下界 $|A|(0)=1$, 嵌入极小圆盘的曲率估计(定理~\ref{thm:2.17})蕴含
$B_s^\Sigma$ 中的总曲率可以随 $s$ 充分大而任意大. 利用这一点以及联系总曲率与
长度的推论~\ref{cor:2.7}, 存在 $R_0=R_0(a_0)$, 使得对任何 $R_3\geq R_0$,
\begin{equation}
a_0R_3\leq\frac{R_3}{4}\int_{B_{R_3/2}^\Sigma}|A|^2
\leq\operatorname{Length}(\partial B_{R_3}^\Sigma).
\label{eq:8.48}
\end{equation}
置 $\beta=2\Omega_1^2\omega$. 推论~\ref{cor:2.7} 与推论~\ref{cor:8.4} 给出
$R_2=R_2(R_0,\beta)$, 使得若 $C\geq R_2$, 则存在
$R_0<R_3<R_2/(2\beta)$ 满足
\begin{equation}
\begin{aligned}
\int_{B_{3R_3}^\Sigma}|A|^2
+\beta^{-10}\int_{B_{2\beta R_3}^\Sigma}|A|^2
&\leq C_2R_3^{-2}\operatorname{Area}(B_{R_3}^\Sigma)\\
&\leq C_2\frac{\operatorname{Length}(\partial B_{R_3}^\Sigma)}{2R_3}.
\end{aligned}
\label{eq:8.49}
\end{equation}
利用 \eqref{eq:8.48}, 取 $n$ 使得
\begin{equation}
a_0R_3\leq4m_1nR_3
<\operatorname{Length}\bigl(\partial B_{R_3}^\Sigma\bigr)
\leq8m_1nR_3,
\label{eq:8.50}
\end{equation}
并固定 $2n$ 条互不相交的曲线
$\tilde\gamma_i\subset\partial B_{R_3}^\Sigma$, 每条长度为 $2m_1R_3$. 定义
内蕴扇区(见图~\ref{fig:8.3})
\begin{equation}
\tilde S_i=\left\{\exp_0(v)\ \middle|\ 0<|v|\leq2\beta R_3
\ \text{且}\ \exp_0(R_3v/|v|)\in\tilde\gamma_i\right\}.
\label{eq:8.51}
\end{equation}
由于各 $\tilde S_i$ 互不相交, 结合 \eqref{eq:8.49} 与 \eqref{eq:8.50} 得
\begin{equation}
\sum_{i=1}^{2n}\left(
 \int_{B_{3R_3}^\Sigma\cap\tilde S_i}|A|^2
 +\beta^{-10}\int_{\tilde S_i}|A|^2\right)\leq4C_2m_1n.
\label{eq:8.52}
\end{equation}
于是, 重新排列各 $\tilde\gamma_i$ 后, 可设对 $1\leq i\leq n$,
\begin{equation}
\int_{B_{3R_3}^\Sigma\cap\tilde S_i}|A|^2
+\beta^{-10}\int_{\tilde S_i}|A|^2\leq4C_2m_1.
\label{eq:8.53}
\end{equation}
由引理~\ref{lem:8.6} 中的 (2), 存在曲线
$\gamma_i\subset\partial B_{2R_3}^\Sigma\cap\tilde S_i$, 长度为 $2m_1R_3$, 使得
若 $y\in S_i=S_{\beta R_3}(\gamma_i)\subset\tilde S_i$, 则
\begin{equation}
B_{\operatorname{dist}_\Sigma(0,y)/2}^\Sigma(y)\subset\tilde S_i.
\label{eq:8.54}
\end{equation}
于是, 由嵌入圆盘的曲率估计(定理~\ref{thm:2.17})与 \eqref{eq:8.53}, 对
$y\in S_i$ 且 $i\leq n$ 有
\begin{equation}
\sup_{\mathcal{B}_{\operatorname{dist}_\Sigma(0,y)/4}(y)}|A|^2
\leq C_3\operatorname{dist}_\Sigma^{-2}(0,y),
\label{eq:8.55}
\end{equation}
其中 $C_3=C_3(\beta,m_1)$. 对 $i\leq n$, \eqref{eq:8.53} 与 Gauss--Bonnet
定理给出
\begin{equation}
\int_{\gamma_i}k_g\leq2\pi+2C_2m_1<(2C_2+2)m_1.
\label{eq:8.56}
\end{equation}
由 \eqref{eq:8.55} 与引理~\ref{lem:8.6} 中的 (4), 存在
$C_4=C_4(C_3)$, 使得对 $i\leq n$,
\begin{equation}
\frac1{2R_3}\leq\min_{\gamma_i}k_g\leq\max_{\gamma_i}k_g\leq\frac{C_4}{R_3}.
\label{eq:8.57}
\end{equation}
反复应用引理~\ref{lem:2.30}(并使用 \eqref{eq:8.55}), 容易看出存在 $\alpha>0$,
使得若 $i_1<i_2\leq n$ 且
\begin{equation}
\operatorname{dist}_{C^1([0,2m_1],\R^3)}
\left(\frac{\gamma_{i_1}}{R_3},\frac{\gamma_{i_2}}{R_3}\right)\leq\alpha,
\label{eq:8.58}
\end{equation}
则
\[
\{z+u(z)N(z)\mid z\in\mathcal{T}_{R_3/4}(S_{i_1})\}
\subset\bigcup_{y\in S_{i_2}}
 B_{\operatorname{dist}_\Sigma(0,y)/4}^\Sigma(y)
\]
对某个 $u\neq0$ 成立, 且
\begin{equation}
|\nabla u|+|A||u|
\leq C_0'\operatorname{dist}_{C^1([0,2m_1],\R^3)}
\left(\frac{\gamma_{i_1}}{R_3},\frac{\gamma_{i_2}}{R_3}\right).
\label{eq:8.59}
\end{equation}
这里 $\operatorname{dist}_{C^1([0,2m_1],\R^3)}
\bigl(\gamma_{i_1}/R_3,\gamma_{i_2}/R_3\bigr)$ 是这两条曲线之间与尺度无关的
$C^1$-距离.

接下来, 我们用紧致性说明当 $n$ 足够大时 \eqref{eq:8.58} 必定成立. 具体地,
由于每个 $\gamma_i/R_3\subset B_2$ 都按弧长参数化在 $[0,2m_1]$ 上, 且由
\eqref{eq:8.57} 有一致的 $C^{1,1}$ 界, 这族映射由 Arzelà--Ascoli 定理给出
紧性. 因此存在 $n_0$, 使得若 $n\geq n_0$, 则对某些 $i_1<i_2\leq n$,
\eqref{eq:8.58} 成立. 特别地, 由 \eqref{eq:8.59} 与引理~\ref{lem:2.29},
当 $n$ 足够大时 $S_{i_1}$ 是 $1/2$-稳定的(现在选取 $a_0,R_0,R_2$). 在旋转
$\R^3$ 之后, 命题~\ref{prop:8.8} 给出 $N$ 值图
$\Sigma_g\subset S_{i_1}$, 其定义在
$D_{2\omega\Omega_1R_3}\setminus D_{2\Omega_1R_3}$ 上, 梯度不超过 $\epsilon$,
$|A|\leq\epsilon/r$, 且
\[
\operatorname{dist}_\Sigma(0,\Sigma_g)\leq8\Omega_1R_3.
\]
按 $r_0/C$ 重新尺度化, 定理即得, 其中
$\bar R=2\Omega_1R_3r_0/C$.
\end{proof}


% ------------------------------------------------------------
%  第 8 章  §2  单侧曲率估计
% ------------------------------------------------------------
\section{单侧曲率估计}
\label{sec:8.2}

理解嵌入极小曲面的关键工具之一, 是 \cite{CM6} 的单侧曲率估计. 接下来的两节
将用这个估计证明两个非常不同的长期猜想: 广义 Nitsche 猜想, 以及嵌入曲面的
Calabi--Yau 猜想.

单侧曲率估计大致说的是: 一个位于平面一侧、但又贴近该平面的嵌入极小圆盘,
其曲率是有界的. 换个说法: 若一个球的中心处曲率很大, 则该极小圆盘向各个
方向传播开来, 因而它不可能被包含在任何经过球心附近的平面的某一侧.

\begin{theorem}[Colding--Minicozzi, \cite{CM6}]
\label{thm:8.9}
存在 $\epsilon_0>0$, 使得下述结论成立. 设 $y\in\R^3$、$r_0>0$, 且
\begin{equation}
\Sigma^2\subset B_{2r_0}(y)\cap\{x_3>x_3(y)\}\subset\R^3
\label{eq:8.60}
\end{equation}
是一个紧的嵌入极小圆盘, 且 $\partial\Sigma\subset\partial B_{2r_0}(y)$. 则对
$B_{r_0}(y)\cap\Sigma$ 的任意连通分支 $\Sigma'$, 只要
$B_{\epsilon_0r_0}(y)\cap\Sigma'\neq\emptyset$, 就有
\begin{equation}
\sup_{\Sigma'}|A_{\Sigma'}|^2\leq r_0^{-2}.
\label{eq:8.61}
\end{equation}
\end{theorem}

单侧曲率估计是在 \cite{CM6} 中证明的, 用到了 \cite{CM3}--\cite{CM6} 的结果.

\mtfigure[0.40]{fig-8-4}{单侧曲率估计.}{fig:8.4}

重新尺度化的悬链面的例子表明, 单连通性与嵌入性这两个假设对单侧曲率估计
都是本质的. 更确切地说, 悬链面的高度关于到旋转轴的距离是对数增长的.
特别地, 悬链面与 $B_{r_0}$ 的交集落在一个厚度约为 $\log r_0$ 的薄板内, 而
\begin{equation}
\frac{\log r_0}{r_0}\longrightarrow0\qquad\text{当 } r_0\to\infty.
\label{eq:8.62}
\end{equation}
因此, 把悬链面按 $\frac1j$ 伸缩后, 我们得到单位球内的一列极小曲面, 它们作为
集合在 $j\to\infty$ 时收敛到 $\{x_3=0\}$. 然而悬链面不是平的, 故这些重新
尺度化的悬链面有 $|A|\to\infty$, 于是当 $j$ 足够大时 \eqref{eq:8.61} 不适用.

我们稍后将回到 \cite{CM24} 中证明的单侧曲率估计的一个内蕴版本.

\mtfigure[0.40]{fig-8-5}{重新尺度化的悬链面.}{fig:8.5}


% ------------------------------------------------------------
%  第 8 章  §3  广义 Nitsche 猜想
% ------------------------------------------------------------
\section{广义 Nitsche 猜想}
\label{sec:8.3}

本节研究一个有限拓扑的真嵌入完全极小曲面 $\Sigma^2\subset\R^3$ 的
端. 特别地, 我们证明广义 Nitsche 猜想; 它是单侧曲率估计的一个应用
(来自 \cite{CM10}). 最初的证明由 Collin \cite{Co} 给出.

给定一个有限拓扑的真嵌入极小曲面, 每个端都有一个代表 $E$, 它是一个
真嵌入的极小环面. 若 $E$ 有有限总曲率, 则它渐近于一个平面或半个
悬链面. 另一方面, 螺旋面提供了唯一一个有限拓扑且总曲率无穷的端的例子.
显然, 螺旋面的端不可能有一个与平面或悬链面的端不相交的代表.

\cite{CM10} 的主要结果是: 任何位于一个充分狭窄的向下倾斜的锥体之上的真嵌入完全极小环面都必有有限总曲率. 这与下文将要描述的 P. Collin
\cite{Co} 的一个结果密切相关.

在 \cite{HoMe2} 中, D. Hoffman 与 W. Meeks 证明了 $\Sigma$ 至多有两个总曲率
无穷的端. 此外他们猜想: 若上述 $\Sigma$ 至少有两个端, 则 $\Sigma$ 必有有限
总曲率(这就是``有限总曲率猜想''; 参见 \cite{Me1}).

若 $\Sigma$ 至少有两个端, 则由 \cite{HoMe2} 的引理 5, 存在一个与 $\Sigma$
不相交的平面端或悬链面端. 因此, 为证明有限总曲率猜想, 只须证明: 位于
半个悬链面下方的真嵌入极小环面端 $E$ 有有限总曲率.

在这个方向上, W. Meeks 与 H. Rosenberg \cite{MeR1} 证明了: 若 $\Sigma$ 至少
有两个端, 则 $\Sigma$ 共形等价于一个去掉了有限多个点的紧 Riemann 面. 利用
这一点, 他们证明了有限总曲率猜想中出现的环面端 $E$ 落在一个半空间内.
事实上, 他们证明了 $E$ 要么渐近于平面(总曲率有限), 要么满足``广义 Nitsche
猜想''的假设.

\begin{conjecture}[广义 Nitsche 猜想, \cite{MeR1}; Collin 定理, \cite{Co}]
\label{conj:8.10}
对 $t\geq0$ 令 $P_t=\{x_3=t\}\subset\R^3$. 设 $E\subset\{x_3\geq0\}$ 是一个
真嵌入的极小环面, $\partial E\subset P_0$, 且 $E$ 与每个 $t>0$ 的
$P_t$ 相交于一条简单闭曲线; 则 $E$ 有有限总曲率.
\end{conjecture}

P. Collin 在 \cite{Co} 中证明了本猜想, 从而结合 \cite{MeR1} 表明: 对至少
有两个端的真嵌入完全极小曲面, 有限拓扑与有限总曲率是等价的. H.
Rosenberg 与 E. Toubiana \cite{RoTo} 构造的一个总曲率无穷的真浸入
极小圆柱面表明, 嵌入性在有限总曲率猜想与广义 Nitsche 猜想中都是必要的
假设. 众所周知的例子还表明, 固有性也是必要的.

设 $x_1,x_2,x_3$ 是 $\R^3$ 上的标准坐标. 给定任意 $\epsilon\in\R$, 记
$C_\epsilon\subset\R^3$ 为锥形区域
\begin{equation}
\mathcal{C}_\epsilon=\left\{x_3>\epsilon\sqrt{x_1^2+x_2^2}\right\}.
\label{eq:8.63}
\end{equation}
按此定义, $C_0$ 是半空间, 而 $C_\epsilon$ 是凸的当且仅当 $\epsilon\geq0$.

\begin{theorem}[Colding--Minicozzi, \cite{CM10}]
\label{thm:8.11}
存在 $\epsilon>0$, 使得任何真嵌入的完全极小环面端
$E\subset\mathcal{C}_{-\epsilon}$ 都有有限总曲率.
\end{theorem}

广义 Nitsche 猜想直接由定理~\ref{thm:8.11} 推出, 因为
$\{x_3\geq0\}=C_0\subset C_{-\epsilon}$. 事实上, 由于悬链面的高度只按对数
增长, 有限总曲率猜想也直接由定理~\ref{thm:8.11} 推出.

单侧曲率估计(定理~\ref{thm:8.9})将是证明的关键.

下面的引理将在证明中被用来在端 $E$ 中寻找``低点''.

\begin{lemma}
\label{lem:8.12}
设 $E$ 如上一定理所述, 且 $\delta>0$, 则存在一列点
$y_j\in E\setminus C_\delta$ 满足 $|y_j|\to\infty$.
\end{lemma}

\begin{proof}
若对某个 $\delta>0$ 这一点不成立, 我们可以重新尺度化使得
$\partial E\subset B_1$ 且
\begin{equation}
E\setminus B_1(0)\subset\mathcal{C}_\delta.
\label{eq:8.64}
\end{equation}
由于 $E$ 是固有的且 $\delta>0$, 存在一条曲线 $\sigma_0\subset E$ 把 $B_1$
与 $\{x_3=200\}$ 相连.

由于 $\delta>0$, 存在 $s_0<0$ 使得对一切 $t\geq1$ 有
$\cosh^{-1}(t)+s_0<\delta t$; 于是
\begin{equation}
\operatorname{Cat}(0,0,s_0)\cap\bar E=\emptyset.
\label{eq:8.65}
\end{equation}
由强极值原理, 不存在第一个 $s>s_0$ 使
$\operatorname{Cat}(0,0,s)\cap\bar E\neq\emptyset$(因为
$\partial E\subset B_1$). 特别地,
\begin{equation}
E\subset\left\{x_1^2+x_2^2<1\right\}.
\label{eq:8.66}
\end{equation}
再次由强极值原理, 不存在第一个 $0<s\leq10$ 使
\begin{equation}
\operatorname{Cat}(s,0,100)\cap\bar E\neq\emptyset.
\label{eq:8.67}
\end{equation}
这就给出矛盾, 因为 $\operatorname{Cat}(10,0,100)$ 在
$\{x_1^2+x_2^2<1\}$ 中把 $B_1$ 与 $\{x_3=200\}$ 隔开.
\end{proof}

在定理的证明中, 除了通常的外蕴球 $B_s(y)$ 之外我们还要用内蕴球. $E$ 中
以点 $y\in E$ 为中心、半径为 $s$ 的内蕴球记作 $\mathcal{B}_s(y)$.

\begin{proof}[定理~\ref{thm:8.11} 的证明]
由初等拓扑学, 存在一条连通曲线 $\beta\subset\mathcal{C}_{-\epsilon}\setminus E$,
其一端在 $\partial E$, 另一端在原点. 由于定理的陈述是伸缩不变的, 我们可以
假设
\begin{equation}
\partial E\cup\beta\subset B_1\cap\mathcal{C}_{-\epsilon}.
\label{eq:8.68}
\end{equation}
引理~\ref{lem:8.12} 给出一列点 $y_j\in E\setminus\mathcal{C}_\epsilon$, 满足
$4<|y_j|\to\infty$. 由 Sard 定理, 我们可以选取一列 $r_j$ 使
\begin{equation}
|y_j|-2\leq6r_j\leq|y_j|-1,
\label{eq:8.69}
\end{equation}
使得 $B_{r_j}(y_j)$ 与 $B_{2r_j}(y_j)$ 横截地与 $E$ 相交.

对每个 $j$, 我们在 $\R^3$ 上定义函数 $f_j$ 为
\begin{equation}
f_j(x)=2\bigl(x\cdot y_j-2|y_j|\bigr)^2-\bigl||y_j|x-2y_j\bigr|^2,
\label{eq:8.70}
\end{equation}
并注意到 $f_j$ 在任何极小曲面上都是超调和的, 因此
\begin{equation}
\Delta_Ef_j\leq0.
\label{eq:8.71}
\end{equation}
记 $\Sigma_j$ 为 $B_{2r_j}(y_j)\cap E$ 中含 $y_j$ 的连通分支. 若 $\Sigma_j$ 不是
圆盘, 则存在一个极小环面 $\Gamma_j$, 其边界分支
$\eta_j^1\subset B_1$ 与 $\eta_j^2\subset B_{2r_j}(y_j)$. 由于 $f_j$ 在
$\eta_j^1$ 与 $\eta_j^2$ 上为正, 极大值原理蕴含 $f_j>0$ 于 $\Gamma_j$ 上.
这是矛盾, 因为 $\eta_j^1$ 与 $\eta_j^2$ 位于 $\{f_j>0\}$ 的不同连通分支中;
故 $\Sigma_j$ 是圆盘.

由构造, 我们有 $x_3(y_j)\leq12\epsilon r_j$, 且
\begin{equation}
B_{2r_j}(y_j)\cap\Sigma_j\subset B_{2r_j}(y_j)\cap\mathcal{C}_{-\epsilon}
\subset\left\{x_3\geq-14\epsilon r_j\right\}.
\label{eq:8.72}
\end{equation}
我们由此断定单侧曲率估计(定理~\ref{thm:8.9})适用. 特别地, 若
$26\epsilon\leq\epsilon_0$, 则
\begin{equation}
\sup_{\Sigma_j^1}|A_{\Sigma_j^1}|^2\leq r_j^{-2},
\label{eq:8.73}
\end{equation}
其中 $\Sigma_j^1$ 是 $B_{r_j}(y_j)\cap\Sigma_j$ 中含 $y_j$ 的连通分支.

曲率估计 \eqref{eq:8.73} 使我们可以把 Harnack 不等式 \cite{CgYa} 应用于
$\mathcal{B}_{r_j}(y_j)\subset\Sigma_j^1$ 上的正调和函数
$x_3+14\epsilon r_j$. 我们得到
\begin{equation}
\sup_{\mathcal{B}_{\frac{3r_j}{4}}(y_j)}(x_3+14\epsilon r_j)
\leq C_h(x_3+14\epsilon r_j)(y_j)\leq26C_h\epsilon r_j,
\label{eq:8.74}
\end{equation}
其中 $C_h$ 来自 Harnack 估计. 应用梯度估计, 我们有
\begin{equation}
\sup_{\mathcal{B}_{\frac{5r_j}{8}}(y_j)}|\nabla x_3|
\leq C_g(26C_h\epsilon),
\label{eq:8.75}
\end{equation}
其中 $C_g$ 来自梯度估计 \cite{CgYa}. 当 $\epsilon>0$ 充分小时, \eqref{eq:8.75}
蕴含 $\mathcal{B}_{\frac{5r_j}{8}}(y_j)$ 是 $x_3=0$ 上一个梯度很小的图.
特别地, 存在一点
\begin{equation}
y_j^1\in\partial B_{|y_j|}\cap
\left\{\bigl(x_1-x_1(y_j)\bigr)^2+\bigl(x_2-x_2(y_j)\bigr)^2=r_j^2/4
\right\}\cap\mathcal{B}_{\frac{5r_j}{8}}(y_j).
\label{eq:8.76}
\end{equation}
由 \eqref{eq:8.74}, 若 $26C_h\epsilon\leq\epsilon_0$, 我们现在可以用 $y_j^1$
代替 $y_j$ 重复前面的论证.

如此迭代至多 $48\pi+1$ 次后, 我们完全绕圆柱一周:
\begin{equation}
\partial B_{|y_j|}\cap
\left\{-\epsilon|y_j|\leq x_3\leq C_h^{48\pi+1}(6\epsilon|y_j|)\right\}.
\label{eq:8.77}
\end{equation}
若 $\epsilon>0$ 充分小, 则单侧曲率估计在整个圆柱上都适用. 因此, 只要
$\epsilon$ 充分小, 迭代便给出一条曲线
\begin{equation}
\gamma_j\subset\partial B_{|y_j|}\cap
\left\{-\epsilon|y_j|\leq x_3\leq C_h^{48\pi+1}(6\epsilon|y_j|)\right\}\cap E,
\label{eq:8.78}
\end{equation}
使得 $\gamma_j$ 在 $\{x_3=0\}$ 上是图形化的. 由于 $\gamma_j$ 是嵌入的, 它要么
无限地螺旋, 要么是闭的并与 $x_3$-轴相绕. 由于 $E$ 是固有的, $\gamma_j$ 是
紧的, 因而是闭的.

令 $E_j\subset E$ 为 $B_{|y_j|}\cap E$ 中含 $\partial E$ 的连通分支. 由极大值
原理, $E_j$ 是环面, 而 $B_{|y_j|}\cap E$ 的其他分支都是圆盘.

我们接下来证明 $\partial E_j=\partial E\cup\sigma_j$, 其中
\begin{equation}
\sigma_j\subset\partial B_{|y_j|}\cap
\left\{-\epsilon|y_j|\leq x_3\leq C_h^{48\pi+1}(6\epsilon|y_j|)\right\}.
\label{eq:8.79}
\end{equation}
若 $\gamma_j\subset\overline{E_j}$, 则由 \eqref{eq:8.78} 这是显然的. 另一方面,
若 $\gamma_j\not\subset\overline{E_j}$, 则 $B_{|y_j|}\cap E$ 中以 $\gamma_j$ 为
界的分支必是一个圆盘 $F_j\subset E$. 这时 \eqref{eq:8.78} 与极大值原理蕴含
\begin{equation}
F_j\subset
\left\{-\epsilon|y_j|\leq x_3\leq C_h^{48\pi+1}(6\epsilon|y_j|)\right\}.
\label{eq:8.80}
\end{equation}
由于 $\partial F_j$ 是图形化的且 $F_j$ 是圆盘, \eqref{eq:8.80} 蕴含 $F_j$ 在
$B_{|y_j|}$ 中把 $\partial\mathcal{C}_{-\epsilon}$ 与
$\{x_3>C_h^{48\pi+1}(6\epsilon|y_j|)\}$ 隔开. 由于 $E$ 是嵌入的且
$F_j\cap\beta=\emptyset$, 我们断定 \eqref{eq:8.79} 成立.

如上所述, \eqref{eq:8.79} 使我们可以应用定理~\ref{thm:8.9} 来证明 $\sigma_j$
是具有一致小梯度的图形, 并且进一步证明在图 $\sigma_j$ 的一个邻域上 $|A|$
有一致的与尺度无关的界. 与上一章中正质量定理的证明一样, 这蕴含
$\sigma_j$ 的总测地曲率有一致的界. 于是, 由于 $E$ 的另一个边界分支
$\partial E$ 与 $j$ 无关, 且每个 $E_j$ 有相同的拓扑, 带边界的 Gauss--Bonnet
定理给出
\begin{equation}
\int_{E_j}|A|^2\leq C_1,
\label{eq:8.81}
\end{equation}
其中 $C_1<\infty$, 且 $E_j^0$ 与 $j$ 无关. 由于当 $j\to\infty$ 时 $E_j$ 穷尽
$E$, \eqref{eq:8.81} 与单调收敛定理蕴含
\begin{equation}
\int_E|A|^2\leq C_1,
\label{eq:8.82}
\end{equation}
证明完毕.
\end{proof}


% ------------------------------------------------------------
%  第 8 章  §4  嵌入曲面的 Calabi-Yau 猜想
% ------------------------------------------------------------
\section{嵌入曲面的 Calabi--Yau 猜想}
\label{sec:8.4}

回顾: $\R^n$ 中的一个浸入子流形称为固有的, 若 $\R^n$ 的任意紧子集的
原像在该曲面中是紧的. 这一性质在极小子流形的理论中起了重要作用, 该领域的
许多经典定理都假定固有性. 容易看出任何紧子流形自动是固有的. 另一方面,
没有理由指望一般的浸入(甚至嵌入)是固有的. 例如, 用极坐标参数化的非紧
曲线
\[
\begin{aligned}
\rho(t)&=\pi+\arctan(t),\\
\theta(t)&=t
\end{aligned}
\]
在半径为 $\pi/2$ 与 $3\pi/2$ 的圆之间无限地螺旋. 然而, 人们长期认为极小
浸入(或嵌入)应当有更好的表现. 这一原理由 Calabi--Yau 猜想刻画, 它们可追溯
到 1960 年代. 在此后的四十年间, 人们对它们做了大量工作. 它们最初的形式由
E. Calabi 于 1965 年在 \cite{Ce} 中给出(另见 S. S. Chern, \cite{Ch1} 第 212 页,
以及 S. T. Yau 1982 年的问题清单 \cite{Ya4}):

\begin{conjecture}
\label{conj:8.13}
``证明 $\R^n$ 中的完全极小超曲面必定是无界的.''
\end{conjecture}

Calabi 接着写道: ``已知 $\R^n$(或任何截面曲率 $\leq0$ 的单连通完全 Riemann
流形)中没有紧的极小子流形. 一个更有雄心的猜想是'':

\begin{conjecture}
\label{conj:8.14}
``$\R^n$ 中的完全[非平]极小超曲面在每个 $(n-2)$ 维平坦子空间中的投影都是
无界的.''
\end{conjecture}

这些猜想的浸入版本后来被证明是假的. 具体地, Jorge 与 Xavier \cite{JXa2} 在
1980 年构造了包含在两个平行超平面之间的非平极小浸入, 给出了更有雄心的
猜想~\ref{conj:8.14} 的浸入版本的反例; 另见 \cite{RoTo}. 另一个重要进展
出现在 1996 年: N. Nadirashvili \cite{Na1} 构造了一个极小圆盘到 $\R^3$ 单位
球内的完全浸入, 表明猜想~\ref{conj:8.13} 对浸入曲面同样失败; 关于圆盘以外
的其他拓扑类型, 参见 \cite{MaMo1}, \cite{LMaMo1}, \cite{LMaMo2}.

由固有性的定义显然可知, $\R^3$ 中的固有的极小曲面必定是无界的, 故
Nadirashvili 的例子不是固有的. 远不那么显然的是: 平面是半空间中唯一
的完全固有的浸入极小曲面. 不过这是 D. Hoffman 与 W. Meeks 的强半空间定理
的一个推论:

\begin{theorem}[Hoffman--Meeks, \cite{HoMe3}]
\label{thm:8.15}
包含于 $\{x_3>0\}\subset\R^3$ 中的完全连通真浸入极小曲面必定是
水平平面 $\{x_3=\text{常数}\}$.
\end{theorem}

在 \cite{CM24} 中, Calabi--Yau 猜想被证明对嵌入曲面成立. 我们将在下文更
精确地描述这一点.

\cite{CM24} 的主要结果是圆盘的固有性的一个有效形式, 它给出一个
弦弧界. 显然内蕴距离大于外蕴距离, 故弦弧界的意义在于反向
不等式, 即用外蕴距离从上方控制内蕴距离. 这由下面的定理实现:

\begin{theorem}[Colding--Minicozzi, \cite{CM24}]
\label{thm:8.16}
存在常数 $C>0$, 使得: 若 $\Sigma\subset\R^3$ 是嵌入极小圆盘,
$B_{2R}^\Sigma=B_{2R}^\Sigma(0)$ 是 $\Sigma\setminus\partial\Sigma$ 中半径为
$2R$ 的内蕴球, 且 $\sup_{B_{r_0}^\Sigma}|A|^2>r_0^{-2}$(其中 $R>r_0$), 则对
$x\in B_R^\Sigma$ 有
\begin{equation}
C\operatorname{dist}_\Sigma(x,0)<|x|+r_0.
\label{eq:8.83}
\end{equation}
\end{theorem}

定理中曲率下界的假设 $\sup_{B_{r_0}^\Sigma}|A|^2>r_0^{-2}$ 是弦弧界
所必需的归一化. 这一点通过对螺旋面重新尺度化并平移即可看出.

完全嵌入极小圆盘的固有性是定理~\ref{thm:8.16} 的直接推论. 具体地, 由
\eqref{eq:8.83}, 当内蕴距离趋于无穷时外蕴距离也趋于无穷. 精确地说, 若
$\Sigma$ 是平的, 因而是平面, 则 $\Sigma$ 显然是固有的; 若它非平, 则对
某个 $r_0>0$ 有 $\sup_{B_{r_0}^\Sigma}|A|^2>r_0^{-2}$, 于是由 \eqref{eq:8.83}
$\Sigma$ 是固有的.

定理~\ref{thm:8.16} 与 \cite{CM6} 的单侧曲率估计(即 \cite{CM6} 中的定理
0.2)一起给出该估计的如下内蕴版本:

\begin{corollary}[Colding--Minicozzi, \cite{CM24}]
\label{cor:8.17}
存在 $\epsilon>0$, 使得若
\begin{equation}
\Sigma\subset\{x_3>0\}\subset\R^3
\label{eq:8.84}
\end{equation}
是嵌入极小圆盘, $B_{2R}^\Sigma(x)\subset\Sigma\setminus\partial\Sigma$, 且
$|x|<\epsilon R$, 则
\begin{equation}
\sup_{B_R^\Sigma(x)}|A_\Sigma|^2\leq R^{-2}.
\label{eq:8.85}
\end{equation}
\end{corollary}

作为这个内蕴单侧曲率估计的推论, Calabi 的第二个、也就是``更有雄心的''猜想
对嵌入极小圆盘也成立.

事实上, \cite{CM24} 证明了 Calabi 的两个猜想以及固有性对有限拓扑的
嵌入曲面也成立. 回顾: 若一个曲面 $\Sigma$ 同胚于一个去掉了有限个点的紧
Riemann 面(这些点称为``穿刺''), 则称 $\Sigma$ 具有有限拓扑. 每个穿刺对应于
$\Sigma$ 的一个端.

下面的半空间定理的推广给出 Calabi 第二个、更有雄心的猜想对有限拓扑嵌入
曲面的版本:

\begin{theorem}[Colding--Minicozzi, \cite{CM24}]
\label{thm:8.18}
平面是半空间中 $\R^3$ 内唯一的具有有限拓扑的完全嵌入极小曲面.
\end{theorem}

\begin{theorem}[Colding--Minicozzi, \cite{CM24}]
\label{thm:8.19}
$\R^3$ 中具有有限拓扑的完全嵌入极小曲面必定是固有的.
\end{theorem}

关于 Riemann 三流形中的固有性已有若干结果. W. Meeks 与 H. Rosenberg
\cite{MeR3} 将其推广, 得到 Riemann 三流形中的局部版本; 他们还把它推广到
$\R^3$ 中具有有限亏格且正单射半径的嵌入极小曲面. 在 \cite{Cb} 中,
B. Coskunuzer 证明了双曲三维空间中面积极小圆盘的固有性, 假设无穷远
边界上至少有一个 $C^1$ 点.

在假定各种曲率界的前提下, 关于固有性与半空间性质有大量工作. Jorge 与
Xavier \cite{JXa1}, \cite{JXa2} 证明了不存在曲率有界的完全浸入极小曲面
包含在 $\bigcap_i\{x_i>0\}$ 中; 后来 Xavier 证明了半空间中唯一的这种曲面是
平面 \cite{Xa}. 最近, G. P. Bessa、Jorge 与 G. Oliveira-Filho \cite{BJO},
以及 H. Rosenberg \cite{Ro} 证明了: 若一个完全嵌入极小曲面曲率有界, 则它
必定是固有的. Meeks 与 Rosenberg \cite{MeR2} 把这个固有性推广到局部
曲率有界且有限拓扑的嵌入极小曲面; 随后 Meeks、J. Perez 与 A. Ros
\cite{MePRs1} 把有限拓扑替换为有限亏格.

受 Nadirashvili 例子的启发, F. Martin 与 S. Morales 在 \cite{MaMo2} 中构造了
一个在(开)单位球中固有的完全有界极小浸入. 也就是说, (开)单位球的紧
子集的原像在该曲面中是紧的, 且该曲面的像在单位球的边界上聚集. 他们在
\cite{MaMo3} 中把这一点推广, 证明 $\R^3$ 的任何凸区域(可能非紧、也可能
不光滑)都容许单位圆盘的一个固有的完全极小浸入. 相关的有趣结果还有若干,
包括 \cite{AFM}, \cite{MaMeNa} 与 \cite{ANa}.

\mtpara{高维情形}
在高维中, 关于超曲面的 Calabi--Yau 猜想所知甚少. 浸入版本的反例都依赖于
Weierstrass 技巧, 而这些技巧不能推广到高维; 目前尚无已知反例反驳 Calabi
的第一个猜想. 然而我们猜想这样的反例是存在的:

\begin{conjecture}
\label{conj:8.20}
对一切 $n>3$, 都存在包含在 $\R^n$ 单位球中的完全浸入极小超曲面.
\end{conjecture}

对于嵌入超曲面, 至少在像 \cite{CM24} 那样假定某些拓扑界的情况下, 我们
猜想:

\begin{conjecture}
\label{conj:8.21}
若 $n>3$ 且 $\Sigma\subset\R^n$ 是可缩的完全嵌入极小超曲面, 则它是无界的.
\end{conjecture}

此外我们猜想固有性也成立:

\begin{conjecture}
\label{conj:8.22}
若 $n>3$ 且 $\Sigma\subset\R^n$ 是可缩的完全嵌入极小超曲面, 则它是真嵌入的.
\end{conjecture}

当 $n=4$ 时, 这些猜想就已经非常有意思了.

高维悬链面表明, 当 $n>3$ 时 Calabi 更有雄心的猜想不成立. 高维悬链面是
嵌入的, 不是平的, 但被包含在两个平行平面之间. $n=3$ 与 $n>3$ 的差别本质
上来自曲面与高维之间调和函数理论的差别. 即: Green 函数在 $\R^2$ 上趋于
无穷, 但在 $\R^k$($k>2$)上保持有界. 这一点之所以相关, 是因为坐标函数是
调和的.

\mtpara{固有性证明的梗概}
嵌入曲面的 Calabi--Yau 猜想的证明超出本书范围, 但我们将结合单侧曲率估计
与我们前面描述的若干关于几乎稳定极小曲面的结果, 给出该证明的一个思路.

\cite{CM24} 中``完全嵌入极小圆盘是固有的''这一证明大致由以下三个主要步骤
组成:

\begin{enumerate}
\item 证明: 若曲面在某个球中是紧的, 则在这个球内我们有好的弦弧界.
\item 证明: 若某个一定大小的球的每个交(与 $\Sigma$ 的交)分支都是紧的
(于是由第一步我们有好的估计), 则与两倍欧氏球的交也都是紧的. 起初可能只
有很差的常数, 但由第一步可得到好的常数.
\item 迭代上述两步.
\end{enumerate}

上面的第一步依赖于 \cite{CM3}--\cite{CM6} 中关于真嵌入极小圆盘的
结果. 这里我们在假定第一步的前提下, 给出第二步证明的梗概.

于是假设: 给定圆盘与所有半径为 $r$ 的欧氏球的交都是紧的, 且有好的
弦弧界. 我们要证明对所有半径为 $2r$ 的欧氏球同样成立.

若不然, 则存在 $B_{2r}\cap\Sigma$ 中同属一个连通分支的两点 $x,y$, 但对某个
充分大的常数 $C$ 有 $\operatorname{dist}_\Sigma(x,y)\geq Cr$. 设 $\gamma$ 是
$B_{2r}\cap\Sigma$ 中连接 $x$ 与 $y$ 的内蕴测地线. 把 $\gamma$ 分成若干段,
我们断定在 $\gamma$ 上必有一对点 $x_0,y_0$ 位于 $B_{2r}$ 中, 它们在内蕴意义
上相距很远, 而在外蕴意义上很近. 我们将从这两点出发, 逐步证明 $x_0$ 与
$y_0$ 在 $B_{2r}\cap\Sigma$ 中不可能连通. 这将给出所需的矛盾.

由假设, $B_r(x_0)\cap\Sigma$ 的每个分支都是紧的, 且由第一步有好的
弦弧界; 因此 $x_0$ 与 $y_0$ 必位于不同的分支中. 于是我们得到
$B_r(x_0)\cap\Sigma$ 的两个紧分支, 它们在中心附近外蕴地接近. 前面已有的
结果(\cite{CM6} 的单侧曲率估计; 见其中的定理 0.2)表明, 这两个分支各自
必有一半带有曲率界. 由于该曲率界是以相联系的球的大小来表述的, 于是可
知这两个分支的某个固定比例部分必是几乎平的——同样是相对于其自身的大小而言.
事实上, 现在容易看出这两个几乎平的区域包含以 $x_0$ 与 $y_0$ 为中心、半径为
$r$ 的某个固定比例的内蕴球. 因此, 我们可以走到这些几乎平的内蕴球的边界,
找到两点 $x_1,y_1$: 每个内蕴球中各一个, 它们外蕴地接近而内蕴地相距很远.

用 $x_1,y_1$ 代替 $x_0,y_0$ 重复上述论证, 得到 $x_2,y_2$. 迭代下去, 我们
得到以 $x_0$ 与 $y_0$ 为中心的大片曲面区域, 上面有先验曲率界. 一旦有了
先验曲率界, 涉及稳定性的改进表明, 甚至这些大片区域也是几乎平的, 因而
不可能在 $B_{2r}$ 中连成一体. 这就是所需的矛盾, 也完成了``嵌入极小圆盘是
固有的''这一证明的第二步梗概.

\mtpara{在有界曲率假设下的弦弧界}
在 Calabi--Yau 猜想证明的第二步梗概中, 单侧曲率估计自然地出现, 用来给出
$\Sigma$ 中彼此靠近但不相交的部分的曲率界. 下面的引理说明如何利用这一
曲率界来证明所期望的``弦弧''型结果(联系外蕴距离与内蕴距离).

设 $\Sigma$ 含 $0$, $R>0$, 记 $\Sigma_{0,R}$ 为 $B_R(0)\cap\Sigma$ 中含 $0$
的连通分支.

\begin{lemma}[Colding--Minicozzi, \cite{CM6} 引理 II.2.1]
\label{lem:8.23}
给定 $R_0$, 存在 $R_1$ 使下述结论成立:

若 $0\in\Sigma\subset B_{R_1}$ 是嵌入极小曲面, $\partial\Sigma\subset
\partial B_{R_1}$, 且
\begin{equation}
\sup_{B_{R_1}^\Sigma}|A|^2\leq4,
\label{eq:8.86}
\end{equation}
则
\begin{equation}
\Sigma_{0,R_0}\subset B_{R_1}^\Sigma.
\label{eq:8.87}
\end{equation}
\end{lemma}

\begin{proof}
设 $\tilde\Sigma$ 是 $\Sigma$ 的万有覆盖, 而
\begin{equation}
\tilde\Pi:\tilde\Sigma\to\Sigma
\label{eq:8.88}
\end{equation}
为覆盖映射. 按照 \cite{CM4} 第 2 节中 $\delta$-稳定的定义(见定义
\ref{def:2.27}), \cite{CM2} 的论证(即关于 $1/2$-稳定曲面的曲率估计; 参见
定理~\ref{thm:2.10})给出常数 $C>10$, 使得若
$B_{CR_0/2}^\Sigma(\tilde z)\subset\tilde\Sigma$ 是 $1/2$-稳定的且
$\tilde\Pi(\tilde z)=z$, 则
\begin{equation}
\tilde\Pi:B_{5R_0}^\Sigma(\tilde z)\to B_{5R_0}^\Sigma(z)
\label{eq:8.89}
\end{equation}
是单射, 且 $B_{5R_0}^\Sigma(z)$ 是一个图, 并满足
\begin{equation}
B_{4R_0}(z)\cap\partial B_{5R_0}^\Sigma(z)=\emptyset.
\label{eq:8.90}
\end{equation}
\cite{CM4} 中的推论 2.13(见推论~\ref{cor:2.31})给出
$\epsilon=\epsilon(CR_0)>0$, 使得若 $|z_1-z_2|<\epsilon$ 且在(互不相交的
球)$B_{CR_0}^\Sigma(z_i)$ 上有 $|A|^2\leq4$, 则每个
\begin{equation}
B_{CR_0/2}^\Sigma(\tilde z_i)\subset\tilde\Sigma
\label{eq:8.91}
\end{equation}
是 $1/2$-稳定的, 其中 $\tilde\Pi(\tilde z_i)=z_i$.

我们断言存在 $n$ 使得
\begin{equation}
\Sigma_{0,R_0}\subset B_{(2n+1)CR_0}^\Sigma.
\label{eq:8.92}
\end{equation}
若不然, 我们得到一条曲线 $\sigma\subset\Sigma_{0,R_0}\subset B_{R_0}$, 它
从 $0$ 通到 $\partial B_{(2n+1)CR_0}^\Sigma$. 对 $i=1,\ldots,n$, 取点
$z_i\in\partial B_{2iCR_0}^\Sigma\cap\sigma$. 于是内蕴球
$B_{CR_0}^\Sigma(z_i)$:

\begin{itemize}
\item 两两不交;
\item 中心位于 $B_{R_0}\subset\R^3$ 中;
\item 满足 $|A|^2\leq4$.
\end{itemize}
特别地, 存在 $i_1$ 与 $i_2$ 使得
\begin{equation}
0<|z_{i_1}-z_{i_2}|<C' R_0n^{-1/3}<\epsilon,
\label{eq:8.93}
\end{equation}
且由 \cite{CM4} 中的推论 2.13(见推论~\ref{cor:2.31}), 每个
$B_{CR_0/2}^\Sigma(\tilde z_{i_j})\subset\tilde\Sigma$ 是 $1/2$-稳定的, 其中
$\tilde\Pi(\tilde z_{i_j})=z_{i_j}$. 由 \eqref{eq:8.90}, 每个
$B_{5R_0}^\Sigma(z_{i_j})$ 是一个图, 并满足
\begin{equation}
B_{4R_0}(z_{i_j})\cap\partial B_{5R_0}^\Sigma(z_{i_j})=\emptyset.
\label{eq:8.94}
\end{equation}
特别地,
\begin{equation}
B_{R_0}\cap\partial B_{5R_0}^\Sigma(z_{i_j})=\emptyset.
\label{eq:8.95}
\end{equation}
这与 $\sigma\subset B_{R_0}$ 连接 $z_{i_j}$ 到 $\partial B_{CR_0}^\Sigma(z_{i_j})$
相矛盾.
\end{proof}


% ------------------------------------------------------------
%  第 8 章  §5  有限亏格的嵌入极小曲面
% ------------------------------------------------------------
\section{有限亏格的嵌入极小曲面}
\label{sec:8.5}

下面我们描述 \cite{CM7} 中关于``非单连通''的有限亏格嵌入极小曲面的两个主要
结构定理. 这些结果, 以及前面描述的关于圆盘的结果, 在极小曲面理论近年的许多
发展中起了关键作用. 我们先讨论平面区域(即亏格为零), 然后转向有限亏格的情形.

\mtpara{一致局部单连通 (ULSC)}
我们说一列曲面是\newterm{ulsc}{一致局部单连通的}{uniformly locally simply connected}(简称 ULSC)的. 一列非单连通的平面区域, 在取子列后自然地分成两种情形, 取决于拓扑是否集中
在点上. 为区分这两种情形, 我们说一列曲面 $\Sigma_i^2\subset\R^3$ 是
\textbf{一致局部单连通的}(uniformly locally simply connected, 简称 ULSC), 若
对每个 $x\in\R^3$ 存在常数 $r_0>0$(依赖于 $x$), 使得对一切 $r\leq r_0$ 及
每个曲面 $\Sigma_i$,
\begin{equation}
B_r(x)\cap\Sigma_i\text{ 的每个连通分支都是一个圆盘.}
\label{eq:8.96}
\end{equation}
例如, 一列重新尺度化的螺旋面是 ULSC 的, 而一列颈部收缩到零的重新尺度化
悬链面则不是. 把 \cite{CM3}--\cite{CM6} 中关于圆盘的结构理论应用于 ULSC
列, 给出这类曲面只有两个局部模型. 即, 在 $\R^3$ 的某个球内局部地, 下列之一
成立:

\begin{itemize}
\item 曲率有界, 且曲面局部是某个平面上的图;
\item 曲率趋于无穷, 且曲面局部是双螺旋阶梯.
\end{itemize}

这两种情形都可以通过取一列螺旋面的重新尺度化来说明: 第一种情形出现在远离
轴的地方, 第二种情形出现在轴上. 若取一列 $\Sigma_i=a_i\Sigma$ 的重新尺度化
螺旋面, 其中 $a_i\to0$, 则曲率沿竖直轴趋于无穷, 而在该轴之外有界. 于是我们
得到:

\begin{itemize}
\item 重新尺度化的螺旋面与远离竖直轴的球的交, 给出平面 $\{x_3=0\}$ 上的一族图.
\item 重新尺度化的螺旋面与以竖直轴为中心的球的交, 给出一个双螺旋阶梯.
\end{itemize}

\mtpara{平面区域}
下面我们描述 \cite{CM7} 中关于``非单连通''的嵌入极小平面区域的两个主要结构
定理.(这些结果的精确陈述及其证明见 \cite{CM7}.) 第一个断言: 任何没有小颈部
的这类曲面都可以由两个取向相反的双螺旋阶梯粘合而成. 注意, 当把两个取向相反
的双螺旋阶梯粘合在一起时, 绕两个轴各一圈后仍停留在同一高度. 第二个给出这类
曲面在有小颈部时的``裤子形曲面''分解: 沿一族短曲线把曲面切开. 切开之后剩下的是
定义在``去掉一个或两个子圆盘的圆盘''之上的图形化片段(一个去掉两个子圆盘的
拓扑圆盘称为裤子形曲面). 这两种结构都出现在称为 Riemann 例子的两参数族极小曲面
中, 作为两个不同的极端.

当该列不再 ULSC 时, 曲面还有别的局部模型. 最简单的例子是一列重新尺度化的
悬链面. 一列重新尺度化的悬链面以重数二收敛到平面. 收敛在光滑拓扑意义下
成立, 但在原点处除外, 那里 $|A|\to\infty$. 这列重新尺度化的悬链面不是 ULSC
的, 因为悬链面上的简单闭测地线(即 $\{x_3=0\}$ 平面中的单位圆)是不可缩的,
而重新尺度化把它收缩到原点. 考虑称为 Riemann 例子的嵌入极小平面区域族, 可以
得到其他类型的曲率爆炸. 模去平移与旋转, 这一族是一个两参数的周期极小曲面
族, 其中参数可以看作颈部的大小以及从一个基本域到下一个基本域的转角. 通过
恰当地选取这两个参数, 可以产生出 Riemann 例子的一列, 它们同时说明两个结构
定理:

\begin{enumerate}
\item 若取一列颈部大小固定而转角趋于 $\frac\pi2$ 的 Riemann 例子, 则转角接近
$\frac\pi2$ 的曲面可以由两个取向相反的双螺旋阶梯粘合而成; 见图~\ref{fig:1.6}.
每个双螺旋阶梯看起来像一个螺旋面. 这一列 Riemann 例子收敛到平行平面的叶状
分解. 收敛在两个螺旋面的轴之外是光滑的(这两条轴正是曲率爆炸的奇异集).
该列是 ULSC 的, 因为颈部大小固定, 因此它说明第一个结构定理.
\item 若取一列颈部大小趋于零的例子, 则得到一列非 ULSC 的例子. 然而这些曲面
可以沿短曲线切开成为一族图形化的裤子形曲面; 见图~\ref{fig:8.7}. 这些短曲线收敛
到点, 而图形化的片段除这些点外收敛到平面, 这说明第二个结构定理.
\end{enumerate}

下面我们描述非单连通嵌入极小平面区域的两个主要结构定理. 每个定理都附带一个
紧致性定理作为推论. 第一个结构定理处理没有小颈部的曲面.

\begin{theorem}[Colding--Minicozzi, \cite{CM7}]
\label{thm:8.24}
任何没有小颈部的非单连通嵌入极小平面区域都可以由两个取向相反的双螺旋阶梯
粘合而成. 此外, 若某点处曲率很大, 则双螺旋阶梯的各层之间的间距很小. 注意,
由于两个双螺旋阶梯取向相反, 绕两个轴各一圈后仍停留在同一高度.
\end{theorem}

下面的紧致性结果, 即\newterm{limit-lamination}{极限层状结构定理}{limit lamination theorem}, 是一个推论:

\begin{corollary}[Colding--Minicozzi, \cite{CM7}]
\label{cor:8.25}
一列 ULSC 的(但非单连通的)嵌入极小曲面, 若曲率趋于无穷, 则有一个子列在
两条直线之外光滑地收敛到平行平面的叶状分解. 这两条直线互不相交, 且与叶状
分解的叶子正交; 这两条直线恰恰是曲率趋于无穷之处.
\end{corollary}

这个推论与圆盘的情形类似, 只是非圆盘时我们得到两条奇异曲线, 而圆盘时只有
一条奇异曲线. 此外, 在这两条直线各自的邻域内, 曲面看起来像绕轴的螺旋面,
且绕两条轴的取向相反. 圆盘情形的精确陈述通常称为极限层状结构定理:

\begin{theorem}[Colding--Minicozzi, \cite{CM6}]
\label{thm:8.26}
设 $\Sigma_i\subset B_{R_i}=B_{R_i}(0)\subset\R^3$ 是一列嵌入极小圆盘,
$\partial\Sigma_i\subset\partial B_{R_i}$, 其中 $R_i\to\infty$. 若
$\sup_{B_1\cap\Sigma_i}|A|^2\to\infty$, 则存在子列 $\Sigma_j$ 以及一条
Lipschitz 曲线 $S:\R\to\R^3$, 使得在旋转 $\R^3$ 之后:

\begin{enumerate}
\item $x_3(\mathcal{S}(t))=t$. (即 $S$ 是 $x_3$-轴上的图.)
\item 每个 $\Sigma_j$ 在 $S$ 之外恰由两个多值图组成(它们一起螺旋).
\item 对每个 $1>\alpha>0$, $\Sigma_j\setminus S$ 在 $C^\alpha$-拓扑下收敛到
$\R^3$ 的叶状分解 $\mathcal{F}=\{x_3=t\}_t$.
\item $\sup_{B_r(\mathcal{S}(t))\cap\Sigma_j}|A|^2\to\infty$ 对一切 $r>0$、$t\in\R$ 成立.
\end{enumerate}
\end{theorem}

在 \cite{Me2} 中, Meeks 利用 \cite{CM3}--\cite{CM6} 与 \cite{MeR2} 证明了极限
层状结构定理中的奇异曲线 $S$ 必定与极限平面正交; 由此 $S$ 是一条直线.

尽管推论~\ref{cor:8.25} 与圆盘的情形相似, 仍值得注意的是圆盘的结果本身并不
足以给出该结论. 也就是说, 尽管 ULSC 列局部由圆盘组成, 圆盘的紧致性结果却是
在半径趋于无穷的全局情形下给出的. 人们可能会错误地以为推论~\ref{cor:8.25}
可以用圆盘的结果加上一个 blow-up 论证来证明. 然而, 可以构造出局部的例子说明
这类论证的困难. \cite{CM18} 构造的局部例子见第 1 章; 其他例子可见
\cite{CaL}, \cite{HoWh2}, \cite{Kh}, \cite{Kl} 与 \cite{MeWe}.

\mtfigure[0.40]{fig-8-6}{极小圆盘的极限层状结构定理. 奇异集 $S$ 与两个多值图
以示意方式画出.}{fig:8.6}

第二个结构定理处理有小颈部的曲面, 并给出裤子形曲面分解.

\begin{theorem}[Colding--Minicozzi, \cite{CM7}]
\label{thm:8.27}
任何有小颈部的非单连通嵌入极小平面区域都可以沿一族短曲线切开. 切开之后
剩下的是定义在``去掉一个或两个子圆盘的圆盘''之上的图形化片段(一个去掉两个
子圆盘的拓扑圆盘称为裤子形曲面). 此外, 若某点处曲率很大, 则所有颈部都非常小.
\end{theorem}

下面的紧致性结果是一个推论:

\begin{corollary}[Colding--Minicozzi, \cite{CM7}]
\label{cor:8.28}
一列非 ULSC 的嵌入极小平面区域, 若曲率趋于无穷, 则有一个子列收敛到一族平坦
的平行平面.
\end{corollary}

\mtparasec{唯一性定理}

\mtparasub{螺旋面的唯一性}
利用单侧曲率估计与 \cite{CM3}--\cite{CM6} 中关于嵌入极小圆盘的结构定理,
Meeks 与 Rosenberg \cite{MeR2} 证明了螺旋面的唯一性:

\begin{theorem}[Meeks--Rosenberg, \cite{MeR2}]
\label{thm:8.29}
平面与螺旋面是 $\R^3$ 中仅有的完全真嵌入单连通极小曲面.
\end{theorem}

定理~\ref{thm:8.29} 中固有性的假设后来被 Colding--Minicozzi 在
\cite{CM24} 中去掉了; 见\secref{sec:8.4}. Catalan 于 1842 年证明: 任何完全可展的极小
曲面要么是平面, 要么是螺旋面.

\mtparasub{亏格一螺旋面: 一般有限亏格的结构结果}
如第 1 章所述, Weber、Hoffman 与 Wolf \cite{HoWW} 最近构造了 $\R^3$ 中一个新的
完全嵌入极小曲面. 他们证明存在一个亏格为一的真嵌入极小曲面, 它只有
一个端, 且该端渐近于螺旋面的端. 这个极小曲面称为亏格一螺旋面. 在伸缩意义下,
亏格一曲面列 $\Sigma_i=a_i\Sigma$($a_i\to0$)像一列曲率趋于无穷的重新尺度化
螺旋面那样收敛到 $\R^3$ 中平行平面的叶状分解. 这一收敛事实上是一个一般性
结果的后果: 上一节中的定理虽是对平面区域陈述的, 对固定亏格的列也以少许改动
的形式成立(精确陈述见 \cite{CM7}).

利用 \cite{CM7} 中关于有限亏格极小曲面的结构结果, J. Bernstein 与 C. Breiner
(参见 Meeks 与 Rosenberg \cite{MeR2}) 证明了任何具有有限亏格且只有一个端的完全
真嵌入极小曲面必渐近于螺旋面:

\begin{theorem}[Bernstein--Breiner, \cite{BB2}; 参见 \cite{MeR2}]
\label{thm:8.30}
任何具有有限亏格且只有一个端的完全真嵌入极小曲面必渐近于螺旋面.
\end{theorem}

\mtfigure[0.42]{fig-8-7}{裤子形曲面分解.}{fig:8.7}

\mtparasub{Riemann 例子的唯一性}
最近, Meeks、Perez 与 Ros 利用 \cite{CM3}--\cite{CM7} 的结构结果以及
\cite{CDM} 中的一个分析结果, 证明了 Riemann 例子的唯一性:

\begin{theorem}[Meeks--Perez--Ros, \cite{MePRs4}]
\label{thm:8.31}
Riemann 例子是仅有的完全真嵌入且有无穷多个端的极小平面区域.
\end{theorem}
